\section{Miscellaneous}

These miscellaneous facts may appear on exams to separate an "A" student from the "B" student.

\begin{itemize}
    \item Turbulence measurements often follow the Rayleigh distribution.
    \item The mean for the Cauchy distribution is undefined.
\end{itemize}

\subsection{Probability of Exceeding the Linear Correlation Coefficient}

The linear correlation coefficient $r$ ranges from -1 to 1. A value of 0 indicates no linear correlation, while a value of 1 indicates a perfect positive linear correlation, and a value of -1 indicates a perfect negative linear correlation. In this class, a table will be given to determine the probability of exceeding a given $r$ for a dataset of size $N$ under the null hypothesis that there is no linear correlation between the variables. Therefore, to find the probability that a dataset indeed has correlation $r$, the probability is \emph{one minus} the value in the table.